\begin{textsample}{POS Dim 4 – 2022\_c6 – Score 3.00 – 2022\_c6/enade\_2022\_tec\_gastr\_1935a4d527.txt}  \label{ex:f4_pos_085}
CAPÍTULO 7 ANÁLISE TÉCNICA DA PROVA Este capítulo tem por objetivo apresentar o desempenho dos estudantes concluintes de Tecnologia em Gastronomia no Enade/2022.
Para isso, foram calculadas as estatísticas básicas da prova em seu todo ( Subseção 7.1.1 ), bem como as estatísticas relacionadas aos Componentes de Formação Geral ( Subseção 7.1.2 ) e de Conhecimento Específico da área ( Subseção 7.1.3 ).
Dadas as suas características, foram analisadas, em separado, as questões objetivas ( Seção 7.2 ) e as questões discursivas ( Seção 7.3 ).
Tomando -se como base as duas questões discursivas do Componente de Formação Geral, nas Subseções 7.3.1.6 a 7.3.1.8, são apresentados os resultados e os comentários sobre a correção das respostas em relação à Língua Portuguesa.
Para as questões objetivas, foram disponibilizados os índices de facilidade e de discriminação Ponto-Bisserial, também em separado, para os Componentes de Formação Geral ( Subseção 7.2.1 ) e de Conhecimento Específico ( Subseção 7.2.2 ).
De cada componente, uma das questões foi escolhida para exemplificar a análise gráfica, relacionando -se as alternativas escolhidas pelos estudantes ( inclusive o gabarito ) com o número de acertos no componente.
No Anexo IV, apresenta -se a íntegra da análise gráfica para todas as questões objetivas.
Para cada uma das questões discursivas, os conteúdos dos tipos mais comuns de respostas dos estudantes são apresentados e comparados com o padrão de resposta esperado ( ver Anexo IX ).
Nas tabelas, constam as seguintes estatísticas das notas¹¹ : média do desempenho na prova, erro-padrão da média, desvio-padrão, nota mínima, mediana e nota máxima para cada um de seus componentes.
Tais estatísticas contemplam o total de estudantes concluintes da área de Tecnologia em Gastronomia inscritos e presentes à prova do Enade/2022, tendo em vista agregações, ou por grandes regiões e o país como um todo, ou por categoria administrativa, organização acadêmica da IES e modalidade de oferta dos cursos.
Em relação aos gráficos de distribuição de notas, o intervalo considerado foi de 10 unidades, aberto à esquerda e fechado à direita, com exceção do primeiro intervalo, [ 0 ; 10 ], fechado em ambos os ¹¹ Uma definição dessas estatísticas pode ser encontrada no Glossário inserido no final deste relatório. extremos.
Para os gráficos de distribuição das notas das questões discursivas, foram consideradas mais duas categorias : questão em branco¹² e nota zero. 7.1 ESTATÍSTICAS BÁSICAS DA PROVA Esta seção apresenta estatísticas selecionadas e histogramas da nota geral ( 7.1.1 ) e de cada componente : Formação Geral ( 7.1.2 ) e Conhecimento Específico ( 7.1.3 ).
São também apresentadas estatísticas selecionadas de subpopulações, caracterizadas por grande região, categoria administrativa, organização acadêmica e modalidade de oferta. 7.1.1 ESTATÍSTICAS BÁSICAS GERAIS Na Tabela 7.1, são apresentadas as estatísticas básicas da prova, por grande região, dos estudantes concluintes de Tecnologia em Gastronomia.
A partir da tabela 7.1, os resultados gerais ( Brasil ) das notas da prova apresentaram média de 47,4, erro-padrão da média de 0,2, desvio-padrão de 14,7, mínima de 0,0, mediana de 47,2 e máxima de 88,4.
Por região, a maior média foi obtida na região Sul, com 51,1 pontos.
O maior erro-padrão de média foi na região Centro-Oeste e o maior desvio-padrão foi na região Sudeste.
A maior nota máxima foi na região Sul ( 88,4 ), e a menor nota máxima foi na região Centro-Oeste ( 83,8 ).
Como mostra o Gráfico 7.1, a moda da distribuição ocorre no intervalo ( 40 ; 50 ], com frequência de aproximadamente 30 \% dos participantes.
A distribuição é crescente até o intervalo modal e diminui conforme o aumento das notas.
Na Tabela 7.2, são apresentadas informações referentes à média da nota final, desagregadas por categoria administrativa, organização acadêmica e modalidade de oferta dos cursos.
Na Tabela 7.2 observa -se que, por categoria administrativa, a maior média foi das instituições públicas, com 50,0 comparado com o 47,1 das instituições privadas.
A mediana dessas categorias foi 49,6 nas instituições públicas e 47,0 nas instituições privadas. 7.1.2 ESTATÍSTICAS BÁSICAS NO COMPONENTE DE FORMAÇÃO GERAL Na Tabela 7.3, são apresentadas as estatísticas básicas em relação ao componente da prova que avalia a Formação Geral dos estudantes concluintes.
A nota de Formação Geral é obtida levando -se em conta os acertos nas oito questões objetivas e a nota média obtida nas Questões Discursivas 1 e 2, sendo a nota de cada questão discursiva obtida levando -se em conta 80 \% da nota relativa ao conteúdo e 20 \% da nota relativa ao desempenho linguístico.
De acordo com a Tabela 7.3, os estudantes do Brasil tiveram um desempenho médio de 46,9, e a média por região foi de 47,9 na região Centro-Oeste ; 47,0 na região Nordeste ; 42,7 na região Norte ; 47,9 na região Sudeste ; e 49,1 na região Sul, sendo essa última a maior média do país, superando a média nacional.
Quanto à variabilidade, em termos nacionais houve 0,3 de erro-padrão da média, e 19,5 de desvio-padrão.
O Centro-Oeste se destaca com o maior erro-padrão da média ( 1,2 ), e o Sul pelo maior desvio-padrão ( 19,9 ).
A mediana nacional foi de 46,9, enquanto dentre as regiões a maior foi a da região Centro-Oeste ( 49,4 ) e a menor a da região Norte ( 41,8 ).
É possível constatar que a maior nota do Brasil foi 97,1 e foi obtida por pelo menos um estudante da região Sudeste.
Quanto à nota máxima das outras regiões : 93,1 na Centro-Oeste ; 92,7 na Nordeste ; 92,5 na Norte ; e 92,7 na Sul.
Nota -se que, no gráfico 7.2, as notas apresentam dispersão um pouco maior que no gráfico 7.1.
Na Tabela 7.4, são apresentadas as informações referentes ao desempenho dos concluintes do Componente de Formação Geral, em diferentes agregações : categoria administrativa, organização acadêmica e modalidade de oferta.
A partir dos dados da Tabela 7.4, observa -se que a maior média foi obtida por estudantes de IES públicas ( 52,1 ) e a menor, pelos de IES privadas ( 46,4 ). 7.1.3 ESTATÍSTICAS BÁSICAS DO COMPONENTE DE CONHECIMENTO ESPECÍFICO Na Tabela 7.5, são apresentadas as estatísticas básicas referentes ao Componente Conhecimento da área de Tecnologia em Gastronomia.
A nota do Componente Específico leva em conta as notas da parte objetiva e a média das notas obtidas nas três questões discursivas da área.
De acordo com a tabela 7.5, a média das notas varia de 43,6 na região Norte a 51,7 na região Sul, com a média nacional em 47,5.
Quanto à variabilidade, em termos nacionais houve 0,3 de erro-padrão da média, e 15,2 de desvio-padrão.
O erro-padrão da média é relativamente baixo em todas as regiões, variando de 0,4 a 0,9.
O desvio-padrão indica a dispersão das notas em relação à média.
Ele varia de 13,8 na região Nordeste a 15,7 na região Sudeste.
A menor nota registrada é 0, enquanto a maior nota é 91.
As notas máximas e mínimas variam entre as regiões, com algumas regiões alcançando pontuações mais altas ou mais baixas.
O comportamento das notas do Componente Específico pode ser observado no Gráfico 7.3, que apresenta um histograma com a distribuição das mesmas : uma distribuição unimodal com a moda no intervalo ( 40 ; 50 ].
Na Tabela 7.6, é apresentada uma comparação dos resultados em relação à categoria administrativa, à organização acadêmica e à modalidade de oferta, levando -se em conta o desempenho de estudantes do Componente Específico da prova.
De acordo com a tabela 7.6, a média das notas varia de 47,3 nas IES privadas a 49,2 nas instituições públicas, sendo que a média geral para as instituições públicas é maior do que para as privadas. 7.2 ANÁLISE DAS QUESTÕES OBJETIVAS Esta seção apresenta estatísticas selecionadas e histogramas das questões objetivas de Formação Geral ( 7.2.1 ) e de Conhecimento Específico ( 7.2.2 ).
São, também, apresentadas e comparadas as médias das subpopulações caracterizadas por grande região.13 7.2.1 COMPONENTE DE FORMAÇÃO GERAL Na Tabela 7.7, são apresentadas as estatísticas básicas relativas às oito questões objetivas do Componente da prova que abrange a Formação Geral dos estudantes.
De acordo com os dados da Tabela 7.7, a média do Brasil foi 49,4.
A menor média foi encontrada na região Norte ( 44,5 ), e a maior, na região Sul ( 52,3 ).
As demais médias foram : 48,9 na região Nordeste, 52,2 na região Centro-Oeste e 50,3 na região Sudeste.
O desvio-padrão do Brasil foi 22,5, sendo o maior desvio-padrão encontrado na região Sul ( 23,0 ), e o menor desvio-padrão encontrado na região Norte ( 21,4 ).
As medianas do Brasil, como um todo, e de quatro regiões ( Centro-Oeste, Nordeste, Sudeste e Sul ) foram iguais a 50,0.
Na região Norte, a mediana foi 37,5.
A nota máxima ( 100,0 ) foi alcançada nas cinco grandes regiões.
As notas mínimas ( 0,0 ) foram iguais para todas as regiões.
Na Tabela 7.8, exposta mais adiante, são apresentados o índice de facilidade e o índice de discriminação ( Ponto-Bisserial ) para cada uma das questões objetivas do Componente de Formação Geral.
De acordo com os dados da Tabela 7.8, as questões objetivas do Componente de Formação Geral, segundo o índice de facilidade, foram assim avaliadas : nenhuma teve o índice de facilidade classificado como “ Muito fácil ” e “ Muito difícil ”.
Duas questões foram classificadas como fácil por terem índices de acertos entre 0,64 e 0,78.
Três questões foram consideradas com índice de facilidade médio, situando -se no intervalo entre 0,49 e 0,58 e três consideradas difíceis, encontrando -se no intervalo entre 0,26 e 0,39.
O índice de facilidade variou de 0,26 a 0,78, e o de discriminação, de 0,39 a 0,55.
Quanto ao índice de discriminação, sete questões obtiveram o índice “ Muito bom ” e apenas um o índice “ Bom ”.
Na Tabela 7.9, é apresentada a distribuição das questões, levando -se em conta, simultaneamente, a classificação dos dois índices. 13 As fórmulas utilizadas para o coeficiente de correlação ponto-bisserial e para o índice de facilidade estão apresentadas no Capítulo 2. desvio-padrão encontrado na região Sul ( 23,0 ), e o menor desvio-padrão encontrado na região Norte ( 21,4 ).
As medianas do Brasil, como um todo, e de quatro regiões ( Centro-Oeste, Nordeste, Sudeste e Sul ) foram iguais a 50,0.
Na região Norte, a mediana foi 37,5.
A nota máxima ( 100,0 ) foi alcançada nas cinco grandes regiões.
As notas mínimas ( 0,0 ) foram iguais para todas as regiões.
Na Tabela 7.8, exposta mais adiante, são apresentados o índice de facilidade e o índice de discriminação ( Ponto-Bisserial ) para cada uma das questões objetivas do Componente de Formação Geral.
De acordo com os dados da Tabela 7.8, as questões objetivas do Componente de Formação Geral, segundo o índice de facilidade, foram assim avaliadas : nenhuma teve o índice de facilidade classificado como “ Muito fácil ” e “ Muito difícil ”.
Duas questões foram classificadas como fácil por terem índices de acertos entre 0,64 e 0,78.
Três questões foram consideradas com índice de facilidade médio, situando -se no intervalo entre 0,49 e 0,58 e três consideradas difíceis, encontrando -se no intervalo entre 0,26 e 0,39.
O índice de facilidade variou de 0,26 a 0,78, e o de discriminação, de 0,39 a 0,55.
Quanto ao índice de discriminação, sete questões obtiveram o índice “ Muito bom ” e apenas um o índice “ Bom ”.
Na Tabela 7.9, é apresentada a distribuição das questões, levando -se em conta, simultaneamente, a classificação dos dois índices.
Os dados da Tabela 7.9 atestam que uma questão teve índice de discriminação “ Bom ” e índice de facilidade “ Difícil ”, duas questões “ Muito bom ” e “ Difícil ”, três questões “ Muito bom ” e “ Médio ” e duas questões “ Muito Bom ” e “ Fácil ”.
No índice de facilidade, as Questões 4,6 e 8 foram classificadas na categoria “ Médio ”, as Questões 1, 5 e 7 foram consideradas difíceis, enquanto as Questões 2 e 3 foram consideradas de nível fácil.
A única questão que teve índice de discriminação “ Bom ” no nível de dificuldade “ Difícil ” foi a Questão 1.
Em particular, a Questão 8 foi a que apresentou o maior poder discriminatório, com índice 0,55, e foi considerada “ Médio ” em termos de facilidade, com uma proporção de 0,50 acertos.
No Gráfico 7.4, para exemplificar, analisa -se o comportamento da Questão 6 de Formação Geral.
No Gráfico 7.4, cada uma das cinco curvas representa o percentual de respostas em determinada alternativa da questão, em função do número de acertos dos estudantes nessa parte da prova ( Formação Geral/Múltipla Escolha ), antes de possíveis eliminações pelo critério Ponto-Bisserial.
Os gráficos relativos às demais questões de Formação Geral constam do Anexo IV. 7.2.2 COMPONENTE ESPECÍFICO Na Tabela 7.10, são apresentadas as estatísticas básicas em relação às questões objetivas do Componente Específico da prova de Tecnologia em Gastronomia por grande região.
Nessa parte da prova, como será discutido ainda nesta subseção, quatro questões não foram usadas no cômputo das notas por terem sido descartadas pelo critério do Ponto-Bisserial.
Assim, as notas foram calculadas com base em 23 das 27 questões objetivas válidas de Componente Específico.
Como se verifica a partir dos dados da Tabela 7.10, a média do Brasil do Componente Específico foi 48,6.
A menor média foi observada na região Norte ( 45,1 ), e a maior, na região Sul ( 52,7 ).
O desvio-padrão de todo o Brasil foi 15,3, sendo o menor desvio-padrão encontrado na região Norte ( 13,7 ), e o maior, na região Sudeste ( 16 ).
A mediana de todo o Brasil foi 50, o mesmo valor da mediana encontrada nas regiões Sudeste e Centro-Oeste.
Nas regiões Norte e Nordeste, as medianas foram, respectivamente, 45,8 e na região Sul 54,2.
A nota máxima da prova foi 91,7, obtida nas questões objetivas do Componente Específico, por, pelo menos, um estudante da região Centro-Oeste e Sudeste.
Em uma região o valor foi 87,5, sendo elas a Norte, Nordeste e Sul.
A nota mínima foi zero, obtida por pelo menos um estudante em cada região e somente 8,3 na região Centro-Oeste.
Na Tabela 7.11, são apresentados os Índices de Facilidade e de Discriminação ( Ponto-Bisserial ) das Questões Objetivas do Componente Específico, para os estudantes de Tecnologia em Gastronomia.
A partir do índice de facilidade obtido, apresentado na Tabela 7.11, pode -se concluir que quase metade das questões objetivas da prova foram consideradas “ Difícil ” : das 27 questões, 12 foram classificadas como “ Difícil ”.
Duas questões foram consideradas “ Muito difícil ”, 4 consideradas como “ Médio ” e 2 como “ Muito Difícil ”.
Somente uma questão foi classificada como “ Muito fácil ”.
Já em relação aos índices de discriminação das questões objetivas do Componente Específico da prova, obtém -se como resultado a seguinte classificação : duas das 27 questões foram consideradas muito boas, enquanto 13 delas tiveram índice de discriminação “ Bom ”.
Entre as demais, oito delas foram classificadas como “ Médio ”, e outras quatro, como “ Fraco ”.
Na Tabela 7.12, é apresentada a distribuição das questões levando -se em conta, simultaneamente, a classificação dos dois índices.
Como revelam os dados das Tabelas 7.11 e 7.12, dentre as questões que alcançaram os maiores índices de discriminação, “ Muito bom ”, uma delas, a Questão 33, teve o maior índice de discriminação ( 0,42 ).
Tal questão foi classificada como “ Médio ”, com índice de facilidade igual a 0,47.
A Questão 29 foi a mais “ Difícil ” entre as 27 questões específicas válidas, com baixo índice de facilidade, apenas 13 \% de acertos.
Essa questão apresentou poder discriminatório fraco, 0,19.
A título de exemplo das análises do comportamento das questões objetivas, no Gráfico 7.5, analisa -se a Questão 35 do Componente Específico.
No Gráfico 7.5, cada uma das cinco curvas representa o percentual de respostas em determinada alternativa da questão, em função do número de acertos dos estudantes nessa parte da prova ( Componente Específico ), antes de possíveis eliminações pelo critério Ponto-Bisserial.
Os gráficos relativos às demais questões do Componente Específico constam no Anexo IV. 7.3 ANÁLISE DAS QUESTÕES DISCURSIVAS Esta seção apresenta estatísticas selecionadas e histogramas das Questões Discursivas de Formação Geral ( 7.3.1 ) e de Componente Específico ( 7.3.2 ).
Expõe e compara, também, as médias de subpopulações, caracterizadas por grande região.
Cumpre notar que nem todas as questões passam pelo crivo dos corretores.
As respostas depois de digitalizadas passam por um pré-processamento de identificação automático.
Algumas destas questões são definidas como um “ branco automático ” ( que não são enviadas para os professores para correção, a não ser que na segunda fase, a de inspeção visual, seja constatada uma resposta ).
O critério usado neste pré-processamento, é verificar para cada questão, a partir da quantidade de PIXELS, o que provavelmente, é uma resposta em branco.
A partir da separação destes documentos que foram considerados em branco, eles são enviados para uma verificação visual, por uma equipe de colaboradores, devidamente capacitados.
Deste modo, quaisquer erros, deste pré-processamento, são encontrados, e as questões preenchidas são enviadas para a correção dos professores.
Eventualmente, algumas respostas em branco não são detectadas nesse pré-processamento automático.
A presença de pequenos riscos, ou manchas, pode fazer o programa não identificar uma questão como um possível branco.
Estas são enviadas para a correção, mas são classificadas como “ branco ” pelos professores, quando da correção. 7.3.1 COMPONENTE DE FORMAÇÃO GERAL As análises dos resultados de desempenho dos estudantes de Tecnologia em Gastronomia, nas duas Questões Discursivas relativas à Formação Geral, encontram -se na Tabela 7.13 e no Gráfico 7.6, cujos dados são analisados de forma sequenciada.
No caso da Tabela 7.13, analisa -se a média das notas de cada Questão Discursiva de Formação Geral que leva em conta as correções realizadas por duas bancas : uma que analisa e pontua aspectos relativos ao conteúdo das respostas, seguindo o padrão de respostas divulgado pelo Inep ( ver Anexo IX ) ; outra que avalia o desempenho linguístico expresso pela redação das respostas, seguindo um padrão de respostas específico da área de Língua Portuguesa, detalhado na seção 7.3.1.7 deste relatório.
A nota de cada questão é composta de 80 \% da nota da correção de conteúdo e 20 \% da nota de desempenho linguístico. outra que avalia o desempenho linguístico expresso pela redação das respostas, seguindo um padrão de respostas específico da área de Língua Portuguesa, detalhado na seção 7.3.1.7 deste relatório.
A nota de cada questão é composta de 80 \% da nota da correção de conteúdo e 20 \% da nota de desempenho linguístico.
Na Tabela 7.13, observa -se que a nota “ média ” nesse conjunto de questões foi inferior à obtida nas objetivas.
Os estudantes, de todo o Brasil, obtiveram em Formação Geral, “ média ” 49,4 nas questões objetivas e 43,2 nas questões discursivas.
Pode -se notar, também, que o desvio-padrão nesse conjunto de questões foi maior do que o obtido nas objetivas : 22,5 nas questões objetivas e 25,2 nas questões discursivas.
A maior “ média ” foi obtida na região Sul ( 44,4 ) e a menor, na região Norte ( 40,1 ).
A “ mediana ” de todo o Brasil, neste componente, foi 47,8.
Nas regiões Centro-Oeste, Nordeste, Norte, Sudeste e Sul as medianas foram 47,8, 49,5, 43,5, 48,3 e 49,3, respectivamente.
As notas “ máxima ” ( 100,0 ) foram obtidas nas regiões Nordeste, e na região Norte ; no Sul a máxima foi de 95,5 ; na região Sudeste 99,0, e na região Centro-Oeste 93,8.
A nota “ mínima ” ( 0,0 ) foi a mesma em todas as regiões do Brasil.
No Gráfico 7.6, está representada a distribuição das notas nas questões discursivas do Componente de Formação Geral.
Como mostra o Gráfico 7.6, a moda da distribuição ocorre no intervalo [ 50 ; 60 ], com frequência de 17,5 \% dos participantes, sendo que, no intervalo [ 0 ; 10 ], inclui -se, além da nota zero, a frequência de estudantes que deixaram esse tipo de questão em branco.
A distribuição é crescente do intervalo [ 10 ; 20 ] até a moda e decrescente a partir da mesma.
Cumpre esclarecer que, tendo em vista que as questões discursivas de Formação Geral são padronizadas, ou seja, constam de todas as provas, os comentários da Banca são os mesmos para todas as carreiras acadêmicas, sendo direcionados a todos os estudantes que participaram do Enade/2022.
A seguir, serão analisados o desempenho linguístico dos estudantes da área de Tecnologia em Gastronomia nas duas questões discursivas de Formação Geral do Enade/2022, apresentando -se os resultados obtidos e os comentários da banca de correção para cada questão. 7.3.1.1 ANÁLISE DE CONTEÚDO DA QUESTÃO DISCURSIVA 1 DO COMPONENTE DE FORMAÇÃO GERAL Os dados de Tecnologia em Gastronomia obtidos a partir das respostas à Questão 1 do Componente de Formação Geral, encontram -se na Tabela 7.14 e no Gráfico 7.7, exposto após a apresentação dos dados da referida tabela.
Verifica -se, a partir dos dados da Tabela 7.14, que na Questão 1 do Componente de Formação Geral – de desempenho abaixo ao da outra questão do mesmo componente –, os estudantes, de todo o Brasil, tiveram “ média ” 40,7.
A maior “ média ” para a Questão 1 foi obtida na região Sudeste ( 41,8 ), e a menor, na região Norte ( 38,0 ).
Quanto à variabilidade das notas, o desvio-padrão de todo o Brasil foi 26,9.
O menor desvio-padrão foi obtido na região Sul ( 26,3 ), e o maior desvio-padrão foi obtido pelas regiões Nordeste e Norte ( 27,5 e 27,6 respectivamente ).
A “ mediana ” do Brasil, como um todo, foi 43,5, sendo igual nas regiões Sudeste e Sul, já nas regiões Centro-Oeste foi de 41,0, Nordeste 45,5 e no Norte 41,0.
A nota “ máxima ” da Questão Discursiva 1 foi 100,0, obtida por, pelo menos, um estudante em todas as regiões do Brasil.
Já a nota “ mínima ” para todas as regiões foi zero.
No Gráfico 7.7, apresenta -se a distribuição das notas na Questão Discursiva 1 do Componente de Formação Geral.
O Gráfico 7.7 mostra que 18 \% dos participantes deixaram em branco a resposta à questão cujos dados são nele retratados.
Desconsiderando -se os que obtiveram nota zero, a moda dessa distribuição ocorre no intervalo ( 50 ; 60 ], com frequência de 16 \% dos participantes.
Relatório Síntese de Área – ENADE/2022 7.3.1.2 COMENTÁRIOS SOBRE A CORREÇÃO DE CONTEÚDO DAS RESPOSTAS À QUESTÃO DISCURSIVA 1 A Questão 1 abordou um tema atual, com dois itens para os estudantes responderem : a ) “ discuta por que os investimentos em educação, ciência, tecnologia e inovação são estratégicos e contribuem para o desenvolvimento de um país ” ; b ) “ explique como o fomento público ao desenvolvimento científico pode atender à justiça social e à equidade, em contextos como o da pandemia de Covid-19 ”.
O tema contempla como os investimentos em educação, ciência, tecnologia e inovação ( ECTI ) cooperam para um país se desenvolver e como o fomento público, vinculado ao desenvolvimento científico, pode atender à justiça social, à equidade e à população em momentos de pandemia.
Esse tema é atual, pois o mundo vivenciou recentemente uma pandemia que, somente no Brasil, deixou mais de 700 mil mortos, sendo exaustivamente comentada na televisão e nos meios digitais.
Logo, almeja -se que os respondentes justifiquem a necessidade de investimentos em ECTI com foco em questões como : investimentos público em P\&D por meio de universidades públicas ; pesquisas realizadas em outros países que possibilitaram o acesso mais célere às vacinas e serviram de modelo para o que foi desenvolvido no Brasil ; pesquisadores de universidade, tanto \textbf{brasileiras} quanto estrangeiras como protagonistas para orientar e informar a respeito da Covid-19 por meio dos canais de comunicação ; desenvolvimento de equipamentos de segurança, como máscara de nariz ou de rosto, respiradores, que foram implementados e viabilizados por pesquisa e extensão em universidades e escolas técnicas.
Esperava -se que os candidatos soubessem o básico a respeito do tema, pois o assunto é bastante atual.
Quanto ao item B, esperava -se que os candidatos abordassem questões como : o papel das universidades e dos centros de pesquisa para a justiça e a equidade ; o investimento público em pesquisa, tecnologia e educação para garantir acesso igualitário da população aos recursos de saúde, podendo fazer um contraponto com os investimentos privados ; a política de distribuição de vacinas gratuitas por meio do Sistema Único de Saúde, que possibilitou o acesso às vacinas no Brasil ; e as questões econômicas mundiais que determinaram uma distribuição desigual de recursos de saúde em escala global, com os países mais ricos retendo a maior parte desses recursos.
Atualmente, é comprovado que os países dominantes são aqueles que se desenvolvem e dominam as ferramentas tecnológicas.
As pesquisas mostram que os países que, mesmo em crise, investiram em ciência, tecnologia e inovação são os que conseguiram se reerguer mais rapidamente e sobressair -se aos demais países em crises pandêmicas.
Portanto, a questão foi elaborada de maneira oportuna e adequada ao universo do estudante universitário brasileiro.
No item A, o estudante poderia focar em diversos argumentos, entre eles : como ocorreu o investimento público em pesquisa e desenvolvimento por meio de universidades públicas e institutos de pesquisa ; como as pesquisas realizadas em outros países possibilitaram acesso mais rápido às vacinas e serviram de modelo para o que foi desenvolvido no Brasil ; como o contingenciamento do Fundo Nacional de Desenvolvimento Científico e Tecnológico ( FNDCT ) prejudicou e ainda prejudica o Brasil, entre diversos outros assuntos.
No item B, era esperado que o estudante explicasse como o fomento público ao desenvolvimento científico pode atender à justiça social e à equidade, em contextos como o da pandemia de Covid-19.
Além disso, deveria abordar a importância das universidades com a justiça social ; como os países que investem nessa área têm grandes vantagens em relação a outros países que não investem ; a distribuição de vacina atende à justiça social por possibilitar que toda a população seja vacinada, independentemente de ser rico ou pobre, por exemplo.
O espaço disponível para a resposta dos dois itens foi adequado e a maioria das respostas ocupou parte das linhas, não utilizando todo o espaço.
Algumas respostas não apresentaram a separação dos itens A e B.
A maioria dos estudantes respondeu de modo semelhante ao que constava do padrão de resposta.
Foi recorrente a exposição do problema sem desenvolvê -lo com clareza, apenas de forma superficial.
A abordagem do candidato sobre o tema tornou -se, em alguns casos, muito genérica.
As respostas consideradas boas atenderam ao que a banca esperava.
Abordaram a importância de investir em CT\&I para desenvolver o país, desdobrando os argumentos e citando como isso poderia ocorrer.
No geral, as notas estavam abaixo do esperado.
Ficou evidente que a maioria dos estudantes possui um conhecimento sobre os dois itens apresentados.
A Questão 1 visou avaliar o desenvolvimento da argumentação de um tema relevante para a sociedade brasileira vigente.
Além disso, houve problemas estruturais nos textos, pois a grande maioria dos estudantes tem uma noção do assunto, mas muitos possuem dificuldade de construir um texto com coesão e coerência, ou seja, não conseguem concatenar os argumentos para defender um ponto de vista.
No geral, em relação à Questão 1, identificou -se que a principal deficiência na formação dos estudantes está relacionada ao aprofundamento do conteúdo e à capacidade de expressar, com clareza e coesão, sua opinião a respeito de determinado assunto. 7.3.1.3 ANÁLISE DE CONTEÚDO DA QUESTÃO DISCURSIVA 2 DO COMPONENTE DE FORMAÇÃO GERAL A Tabela 7.15 apresenta dados estatísticos a respeito das notas de conteúdo da Questão Discursiva 2 do Componente de Formação Geral, organizados por grande região de funcionamento dos cursos, no país.
A Tabela 7.15 mostra que o desempenho médio dos estudantes, na Questão Discursiva 2 ( média 45,7 ), foi acima do obtido na Questão Discursiva 1 ( média 40,7 ).
A região Norte foi aquela cuja média foi menor ( 42,1 ), e a maior média foi na região Sul ( 47 ).
Quanto à variabilidade das notas, o desvio-padrão de todo o Brasil foi 29,1, inferior ao obtido na Questão Discursiva 1 ( 26,9 ).
O maior desvio nessa questão foi obtido na região Centro-Oeste ( 30,2 ), enquanto o menor foi obtido na região Sul ( 28,6 ).
A mediana de todo o Brasil foi de 53,0, a mesma da região Sudeste, sendo a região Centro-Oeste de 51,5, Nordeste 53,5, Norte 48,5 e Sul de 53,5.
A nota máxima da Questão Discursiva 2 foi 100,0, obtida por pelo menos um estudante nas quatro regiões.
Já a nota mínima para todas as regiões foi zero.
No Gráfico 7.8, mostra -se a distribuição das notas na Questão Discursiva 2 do Componente de Formação Geral.
Observa -se, a partir dos dados apresentados no Gráfico 7.8, que 20,0 \% dos estudantes deixaram a questão em branco.
O percentual total de zeros chega a 22,5 \%, se agregado à frequência daqueles que, tendo respondido, receberam nota zero.
Nota -se, ainda, que as notas ficaram mais dispersas em comparação às da Questão Discursiva de número 1, o que pode ser constatado, também, pela comparação do desvio-padrão das notas da Questão Discursiva 2 ( 29,1 ) e o das notas da Questão Discursiva 1 ( 26,9 ).
Desconsiderando -se os que obtiveram nota zero, a moda dessa distribuição ocorre no intervalo ( 60 ; 70 ], com frequência de 16,3 \% dos participantes. 7.3.1.4 COMENTÁRIOS SOBRE A CORREÇÃO DE CONTEÚDO DAS RESPOSTAS À QUESTÃO DISCURSIVA 2 A Questão Discursiva 2 abordou um tema importante, com o objetivo de analisar conhecimentos gerais do aluno.
O tema contemplou o conhecimento geral do patriarcado estrutural ( dominação masculina ), o qual é recorrente nos noticiários, e um infortúnio na vida dos brasileiros, pois trata -se de uma questão de violência contra as mulheres.
É recorrente assistir ou ler notícias sobre o domínio masculino na sociedade brasileira.
Questões como a participação da mulher na política ; a paridade de salários entre profissionais de sexos diferentes ; a desigualdade social e o tratamento hostil que é dado às mulheres, em diversas situações, tanto profissionais quanto familiares, devem ser de conhecimento geral do estudante.
Não se objetiva tratar questões de ordem específica para a solução deste problema, mas, apenas, discorrer e argumentar sobre ele.
Portanto, a questão foi elaborada de maneira oportuna e adequada para o universo do estudante universitário brasileiro.
A Questão está dividida em dois itens : A e B.
No item A, esperava -se que o estudante abordasse o patriarcado – dominação masculina que gera a divisão sexual do trabalho, a desigualdade econômica e política, entre homens e mulheres, a sub-representação das mulheres na política e em cargos de poder.
Esperava -se que o estudante ressaltasse, ainda, questões culturais, como a educação sexista e a mídia, que coisificam a imagem das mulheres e naturalizam a desigualdade entre os sexos, de forma a considerar os homens como dominantes, enquanto as mulheres são historicamente tidas como submissas e incapazes.
No item B, era esperado que o estudante apresentasse a importância do papel do Estado na implementação de políticas públicas voltadas à proteção da mulher, tais como na área da saúde, na educação não sexista, na assistência social, na segurança pública – delegacias femininas, defensoria pública especializada, casa de abrigo, juizado especial etc., bem como ONGs que atuam no combate à violência contra a mulher e ao feminicídio.
O espaço disponível para a resposta dos dois itens foi adequado e grande parte das respostas ocupou parte das linhas, não contemplando o espaço completo.
Com relação à linguagem utilizada pelos estudantes, as respostas estavam dentro do esperado.
Alguns estudantes utilizavam termos do jargão do cotidiano.
No item B, o estudante deveria propor duas ações do Estado para o enfrentamento da violência contra as mulheres.
Os alunos deveriam apresentar propostas e desenvolvê -las.
Muitos estudantes responderam apenas por meio da apresentação de ações do Estado, mas sem desenvolvê -las ou argumentá -las.
Em geral, os alunos apenas expuseram ações – algumas vezes, uma, duas e, até mesmo, três ações –, mas sem fazer relações entre elas ou sem desenvolvê -las.
As respostas dos estudantes demonstraram que a maioria apresenta um conhecimento geral a respeito do patriarcado estrutural e da trágica desigualdade que esta impõe dos homens sobre as mulheres, bem como de possíveis ações do Estado para a proteção delas.
A Questão visou avaliar o desenvolvimento da argumentação de um tema relevante para a sociedade brasileira vigente. 7.3.1.5 CONSIDERAÇÕES FINAIS SOBRE A CORREÇÃO DE CONTEÚDO DAS QUESTÕES DISCURSIVAS DO COMPONENTE DE FORMAÇÃO GERAL As Questões Discursivas da área de Formação Geral abordaram temas atuais e proporcionaram a avaliação de competências importantes dentre as previstas pelas suas Diretrizes Curriculares.
A análise das respostas revelou, principalmente, dificuldade dos estudantes em relação à seleção de argumentos sólidos para discorrer sobre o tema – que pode ser fruto de desinteresse ou de desconhecimento mais profundo do contexto.
Esses fatos parecem ter sido os principais responsáveis pelo resultado aquém do esperado em termos de avaliação quantitativa.
Mesmo os estudantes que responderam em conformidade com o padrão demonstraram um universo de argumentação muito limitado, pobre em vocabulário, com o recurso frequente de procurar nos textos introdutórios os elementos para a construção de suas frases.
Encadear e relacionar argumentos foi uma deficiência percebida.
As razões eram apresentadas de forma confusa, sem fluência e sem constituir unidade, utilizando termos genéricos e vagos. 7.3.1.6 ANÁLISE DE LÍNGUA PORTUGUESA DAS QUESTÕES DISCURSIVAS DO COMPONENTE DE FORMAÇÃO GERAL Os dados de Tecnologia em Gastronomia obtidos a partir das respostas às Questões Discursivas do Componente de Formação Geral, no que tange à Língua Portuguesa, são apresentados na Tabela 7.16.
Os dados da Tabela 7.16 mostram que, em relação ao quesito Língua Portuguesa, os estudantes, de todo o Brasil, obtiveram média 52,9.
A maior média em relação à Língua Portuguesa foi obtida na região Sul ( 54,7 ), e a menor, na região Norte ( 49,7 ).
Quanto à variabilidade das notas, o desvio-padrão de todo o Brasil foi 29,4.
O menor desvio-padrão foi obtido na região Sudeste ( 28,3 ), e o maior desvio-padrão foi obtido na região Centro-Oeste ( 31,0 ).
A mediana das notas de Língua Portuguesa foi 57,5 para o Brasil.
As regiões Centro-Oeste, Nordeste e Norte obtiveram 55,0, e as regiões Sudeste e Sul 65,0.
A nota máxima para todo o Brasil foi 100,0, com, pelo menos, um estudante tirando essa nota em todas as regiões.
A nota mínima foi zero em todas as regiões do país.
No Gráfico 7.9, mostra -se a distribuição das notas de Língua Portuguesa do Componente de Formação Geral.
Com base nos dados expostos no Gráfico 7.9, observa -se que a maior frequência ( 20,0 \% ) corresponde a dos estudantes que obtiveram nota no intervalo ( 60 ; 70 ].
Destacam -se, também, os estudantes que deixaram a questão em branco, representando cerca de 17 \% do total de participantes. 7.3.1.7 COMENTÁRIOS SOBRE A CORREÇÃO DE LÍNGUA PORTUGUESA DAS RESPOSTAS ÀS QUESTÕES DISCURSIVAS Os comentários feitos neste item a respeito da correção de Língua Portuguesa das respostas emitidas para as questões discursivas do Componente de Formação Geral estão agrupados em três eixos : avaliação de desempenho linguístico ; padrão de resposta utilizado na avaliação e análise das respostas dadas pelos estudantes.
De acordo com a Portaria nº 298 do MEC/INEP, de 30 de junho de 2022, o Enade deve avaliar se o estudante desenvolveu, no processo de formação, entre outras, competências para [ … ] V - Compreender as linguagens e suas respectivas variações ; VI - Ler, produzir e interpretar textos com clareza e com coerência ; VII - Identificar, analisar e interpretar diferentes representações verbais, não verbais, gráficas e numéricas de fenômenos diversos ou de um mesmo significado ; VIII - Formular e articular argumentos e contra-argumentos consistentes em situações sociocomunicativas ; IX - Identificar, compreender e analisar situações-problema utilizando pensamento holístico e sistêmico ao se abordar a complexidade da realidade. ( INEP, 2022d ) Essas competências são essenciais à formação de nível superior, porque dizem respeito ao domínio da linguagem, tanto no processo de leitura quanto na articulação e na organização de ideias para a produção textual.
Dessa forma, as duas questões discursivas de formação geral oferecem oportunidade para avaliar se o participante consegue desenvolver um tema, recorrendo a argumentos que explicitem sua opinião e elaborando propostas de intervenção social adequadas a um exercício consciente da cidadania.
É, portanto, necessário identificar, nas respostas discursivas de Formação Geral, as diferentes habilidades necessárias ao bom desempenho linguístico na situação de formalidade exigida : a organização das ideias, tendo em vista a construção de um texto coerente e estruturado ; o domínio das convenções ortográficas e das convenções morfossintáticas que constituem o arcabouço da gramática da norma-padrão da Língua Portuguesa ; e o domínio de um vocabulário amplo e diversificado, que permita uma comunicação clara e eficiente.
A Questão Discursiva 1 desenvolveu o tema de investimentos públicos em educação, ciência, tecnologia e inovação, que são estratégicos e contribuem para o desenvolvimento científico de um país, podendo abordar essa questão tanto no Brasil quanto em outros países.
O enunciado estava claro ao encaminhar a solicitação de resposta, evidenciando a exigência de elaboração de um texto que discutisse e explicasse por que os investimentos públicos em educação, ciência, tecnologia e inovação são estratégicos e contribuem para o desenvolvimento científico de um país e como o fomento público ao A análise do desempenho linguístico esperado de um estudante que está concluindo seu curso superior apoiou -se no princípio de que ele deveria utilizar conhecimentos sobre os temas e estruturar seus textos de acordo com as características do registro formal adequado à situação comunicativa – avaliação de conhecimentos ao final de um curso superior.
Essa configuração determina exigências quanto à adequação da seleção vocabular, ao desenvolvimento do conteúdo, à estruturação sintática dos períodos, à organização lógica das ideias, à utilização de procedimentos de encadeamento textual e de referenciação, à obediência às exigências morfossintáticas próprias da modalidade escrita da norma-padrão, ao respeito às regras ortográficas de acentuação gráfica e de grafia das palavras. a ) O padrão de resposta utilizado na avaliação Em relação a este eixo de análise, foram considerados, no padrão de resposta, as seguintes competências : ( 1 ) estruturação textual ; ( 2 ) respeito às convenções ortográficas da norma-padrão da Língua Portuguesa ; ( 3 ) domínio dos diferentes aspectos morfossintáticos próprios da modalidade escrita formal da norma-padrão da Língua Portuguesa ; e ( 4 ) seleção vocabular adequada à modalidade escrita desenvolvimento científico pode atender à justiça social e à equidade, em contextos como o da pandemia de Covid-19.
Os três textos motivadores favoreceram a formulação e a articulação de argumentos e contra-argumentos consistentes, em defesa do ponto de vista a ser defendido pelo participante.
A Questão Discursiva 2 desenvolveu o conceito de patriarcado, tomando como base um texto de Biroli.
Com base nesse texto, o comando solicitou do participante : “ discorra sobre a relação entre o patriarcado estrutural e a desigualdade entre homens e mulheres ” e “ proponha duas ações do Estado para o enfrentamento dos diferentes tipos de violência contra a mulher ”.
O comando “ discorra ” demandava a elaboração de um texto de base dissertativa, com desenvolvimento de argumentação sobre o tema.
Em relação à segunda parte, o comando exigia a elaboração de duas propostas de ações, que deveriam conter informações sobre o agente da intervenção ( Estado ) para o enfrentamento da situação apresentada.
As duas questões, portanto, assemelhavam -se, pois solicitavam um texto dissertativo ao apresentarem comandos para “ discutir ” e “ discorrer ”.
A Questão Discursiva 2, entretanto, solicitava que o aluno apresentasse propostas de ações para o enfrentamento do problema destacado.
Esses dois estilos textuais ofereceram oportunidade de avaliar se os participantes estavam habilitados a desenvolverem textos com esses dois distintos objetivos comunicativos, o que seria esperado ao final do curso de graduação e necessário à sua vida profissional.
A análise do desempenho linguístico esperado de um estudante que está concluindo seu curso superior apoiou -se no princípio de que ele deveria utilizar conhecimentos sobre os temas e estruturar seus textos de acordo com as características do registro formal adequado à situação comunicativa – avaliação de conhecimentos ao final de um curso superior.
Essa configuração determina exigências quanto à adequação da seleção vocabular, ao desenvolvimento do conteúdo, à estruturação sintática dos períodos, à organização lógica das ideias, à utilização de procedimentos de encadeamento textual e de referenciação, à obediência às exigências morfossintáticas próprias da modalidade escrita da norma-padrão, ao respeito às regras ortográficas de acentuação gráfica e de grafia das palavras. a ) O padrão de resposta utilizado na avaliação Em relação a este eixo de análise, foram considerados, no padrão de resposta, as seguintes competências : ( 1 ) estruturação textual ; ( 2 ) respeito às convenções ortográficas da norma-padrão da Língua Portuguesa ; ( 3 ) domínio dos diferentes aspectos morfossintáticos próprios da modalidade escrita formal da norma-padrão da Língua Portuguesa ; e ( 4 ) seleção vocabular adequada à modalidade escrita formal da Língua Portuguesa, exigida pela situação comunicativa.
A seguir, são apresentados os aspectos avaliados associados a cada uma destas competências. ( 1 ) Estruturação textual condizente com o gênero textual solicitado e o modo de organização textual adequado ao gênero ( argumentativo ou expositivo ).
Esta competência envolve : - estruturação sintática condizente com o padrão da modalidade escrita formal da Língua Portuguesa, de modo a garantir a clareza e a coerência necessárias ; - distribuição do conteúdo do texto em parágrafos, de modo a garantir a sua organização temática ; - distribuição das ideias em períodos que mantenham uma estruturação sintática condizente com a modalidade formal da língua ( evitando -se frases fragmentadas, sequência justaposta de ideias sem encaixamentos sintáticos e presença de oração subordinada sem oração principal que reflitam hábitos da oralidade e comprometam a estrutura lógico-gramatical do texto ) ; - utilização de operadores discursivos que contribuam para a progressão temática do texto, estabelecendo relações lógicas entre as ideias apresentadas, tanto do ponto de vista intrafrasal, como interfrasal : preposições, conjunções, pronomes relativos, advérbios e locuções adverbiais de tempo e lugar ; - utilização de procedimentos de referenciação lexical e pronominal que permitam a retomada de referentes textuais, como pronomes, advérbios e artigos ; além de elementos lexicais, como sinônimos, hiperônimos, hipônimos, expressões metafóricas e expressões resumitivas ) ; - emprego adequado dos pronomes relativos “ cujo(a) ” e “ onde ”, responsáveis por encaixamentos sintáticos no período composto ; - utilização adequada de sinais de pontuação que contribuam para a organização lógica da frase e do texto ; - inteligibilidade relacionada ao atendimento das exigências de estruturação textual.
A expectativa, portanto, era de que o participante recorresse a procedimentos linguístico-discursivos para organizar seu texto, permitindo o encadeamento lógico entre suas partes de forma a garantir a progressão e a coerência textuais, com estruturação que não reproduzisse hábitos da oralidade e sim atendesse à norma-padrão da Língua Portuguesa. ( 2 ) Respeito às convenções ortográficas da norma-padrão da Língua Portuguesa.
Esta competência envolve o respeito às convenções estabelecidas pela legislação em vigor e consubstanciadas no Vocabulário Ortográfico da Língua Portuguesa, editado pela Academia Brasileira de Letras e avaliam -se : - domínio das regras de acentuação gráfica ; - domínio da grafia padrão das palavras ( com ausência de abreviaturas próprias da linguagem da internet ).
Assim, esperava -se que o participante : grafasse corretamente as palavras ; respeitasse as regras de acentuação gráfica ; empregasse maiúsculas em início de frase, em nomes próprios de pessoas, lugares ou instituições ; grafasse as siglas com letras maiúsculas ou apenas com a primeira letra maiúscula quando formarem uma palavra ; evitasse abreviações como p/, vc, tb, pq, tá, né, usadas muitas vezes em escrita informal e na internet ; obedecesse às regras de separação de sílabas no final da linha ; atendesse aos padrões de segmentação vocabular da norma-padrão. ( 3 ) Domínio dos diferentes aspectos morfossintáticos próprios da modalidade escrita formal da norma-padrão da Língua Portuguesa.
Esta competência envolve o domínio das convenções próprias da norma-padrão da Língua Portuguesa relativas aos seguintes aspectos : concordância nominal ; concordância verbal ; regência nominal ; regência verbal ; flexão nominal ; flexão verbal ; correlação entre os tempos verbais ; colocação pronominal. ( 4 ) Assim, esperava -se que o participante : - flexionasse o verbo para estabelecer concordância de número com o sujeito da frase ; - flexionasse o artigo, o adjetivo e o pronome para concordar em número e em gênero com o substantivo a que cada um se refere ; - observasse a regência nominal e a verbal, utilizando a preposição adequada depois de um substantivo, um verbo ou um adjetivo ; - empregasse adequadamente o acento grave indicador da crase que expresse a contração entre uma preposição e um artigo, de acordo com as exigências de regência nominal ou verbal ; - obedecesse às regras de colocação pronominal ( próclise e ênclise ), distintas dos hábitos da oralidade ou da escrita informal ; - flexionasse adequadamente verbos, substantivos, adjetivos e pronomes, no que diz respeito à expressão das categorias gramaticais ; - flexionasse os verbos para expressar a correlação de modo e de tempo nas estruturas subordinadas.
Com base nesses critérios, alguns desvios ortográficos foram considerados como desvios de caráter morfossintático porque interferem na classificação dos vocábulos de base morfológica e funcional : - eliminação da marca de infinitivo ( -r- ) e substituição por acento agudo ou ausência total de marca do infinitivo, causando ambiguidade entre distintas classes de palavras ; - confusão entre as terminações “ ão ” e “ am ”, utilizadas para marcar diferentes tempos verbais ; - emprego indevido da forma verbal de"haver"\&lt ; “ há”\&gt ; no lugar da preposição ou do artigo \&lt ; “ a”\&gt ;, e vice-versa ; - emprego de hífen para separar pronome átono – tanto uso indevido quanto omissão ( ex. : “ esperasse ”, em lugar de “ espera -se ” ; “ falar -mos ”, no lugar de “ falarmos ” ) ; - emprego de “ esta ” no lugar de “ está ”, provocando ambiguidade devido à confluência gráfica entre o pronome demonstrativo e a forma do verbo"estar", pela ausência da acentuação gráfica ; - emprego da conjunção aditiva"e"no lugar da forma verbal “ é ”, provocando ambiguidade devido à confluência gráfica por ausência da acentuação gráfica ; - emprego do advérbio “ mais ” no lugar da conjunção adversativa “ mas ”, reproduzindo confluência fônica da oralidade ; - registro da terceira pessoa do plural dos verbos “ ter ” e “ vir ” sem o acento circunflexo, causando desvio de concordância verbal ; - emprego inadequado do relativo “ onde ” deslocado do seu sentido locativo, reproduzindo tendência da oralidade ; - emprego do pronome relativo “ que ” sem a necessária preposição ( desvio de regência ) ; - ausência do acento grave obrigatório para marcar a crase. ( 5 ) Seleção vocabular adequada à modalidade escrita formal da Língua Portuguesa, exigida pela situação comunicativa.
Esta competência envolve : - precisão na seleção/utilização do vocabulário relacionado à temática solicitada pela questão ; - diversificação e adequação vocabular ; - coerência na combinação vocabular ; - ausência de marcas da oralidade, como termos de sentido muito genérico ( “ coisa ”, “ negócio ”, “ você ” ) e termos de registros mais informais ( como gírias, jargões, frases feitas, ditados populares, termos regionais ).
Assim, esperava -se que o participante : respeitasse a adequação vocabular, não empregando gírias ou expressões coloquiais ; apresentasse repertório vocabular diversificado e preciso ; evitasse repetição desnecessária de palavras ; utilizasse um vocabulário mais formal, como solicitado por um texto dissertativo.
Essas quatro competências foram reunidas em três itens de avaliação, atendendo aos aspectos explicitados anteriormente : I - Aspectos Ortográficos ( 30 \% ) : domínio das convenções ortográficas : grafia de vogais e consoantes, uso de maiúsculas e minúsculas, emprego do hífen e da acentuação gráfica, segmentação de palavras ; II - Aspectos Textuais ( 20 \% ) : domínio dos procedimentos de estruturação textual : organização interna dos períodos, emprego de conectores para a articulação lógica entre os períodos e entre os parágrafos, emprego de marcas de referenciação lexical e pronominal ; utilização dos sinais de pontuação que contribuem para a organização lógica da frase ; III - Aspectos Morfossintáticos e Vocabulares ( 50 \% ) : domínio das regras de caráter morfossintático estabelecidas como modelares do ponto de vista da modalidade escrita formal da norma-padrão da Língua Portuguesa : concordância nominal e verbal, regência nominal e verbal, colocação pronominal, flexão nominal e verbal, correlação entre tempos e modos verbais, ausência de marcas de oralidade.
Outra competência combinada a esta na análise das respostas é a seleção vocabular adequada à modalidade escrita formal da Língua Portuguesa, tendo em vista a intersecção entre as duas do ponto de vista das exigências do registro formal da modalidade escrita da norma-padrão. b ) Análise das respostas dadas pelos estudantes A Questão 1 propiciou a oportunidade de maior desenvolvimento do tema solicitado, gerando consequências na estruturação textual das respostas, já que os textos foram mais expositivos, mais elaborados e fluentes, inclusive, um dos textos era uma tabela.
Já a Questão 2 teve um comportamento distinto, devido ao seu caráter mais opinativo, o texto foi mais curto e o vocabulário mais limitado.
Observaram -se, nas duas questões, alguns casos de comprometimento de autoria, com textos muito próximos a paráfrases dos textos motivadores e sem marcas de repertório sociocultural produtivo.
Apresentam -se, a seguir, algumas \textbf{observações} sobre o desempenho dos participantes : I - Aspectos Ortográficos : o desempenho dos participantes revelou uma diferença muito grande nos dois aspectos analisados, como já foi observado nos relatórios dos anos anteriores : baixo índice de desvios da grafia padrão e grande índice de desvios de acentuação.
Em vários casos, ocorre ausência quase completa de acentuação gráfica.
Vale ressaltar que muitos participantes apresentaram uma caligrafia quase ininteligível.
Ainda nesse aspecto, ocorreu a omissão inadequada de cedilha, de til e do traço do T em algumas respostas.
Houve, inclusive, respostas com ausência total de cedilha ou til.
Os casos mais sistemáticos de eliminação do acento indicador da sílaba tônica foram : palavras proparoxítonas ( “ politica ”, “ pratica ”, “ publica ”,"policia") ; palavras paroxítonas terminadas em ditongo crescente ( “ necessario ”,"historia","propria", “ violencia ” ) ; palavras oxítonas ( “ tambem ”,"populacao", “ pais ” ).
Por outro lado, destaca -se o uso indevido do acento gráfico em palavras como “ investímento ” e “ pátriarcado ”, por exemplo.
Quanto aos desvios de grafia, observam -se os seguintes exemplos : emprego inadequado de maiúscula : “ Leis ” ; ou de minúscula :"constituição"“ Ministério publico ” ; confusão entre"ss","sc","c"e"ç": “ nessessaria", “ agreção ”, “ conciencia ”, “ comecam ”, “ cocientização ”, “ pacos ” ; eliminação de grafemas : “ mostando ”,"deveriamo"; - desvios de segmentação : “ des de", “ des igualdade ”, “ pre conceito"; troca de grafema : “ órgam ” ( por “ órgão ” ), “ escravido ” ( por “ escravidão ” ).
II - Aspectos Textuais : a estruturação textual foi a competência que apresentou o maior número de desvios, sendo responsável, em alguns textos, pela falta de clareza e pela desorganização.
No conjunto dos textos analisados, observaram -se diferentes níveis de domínio dos recursos de estruturação e organização textual.
No entanto, os participantes podem ser agrupados em três faixas descritas a seguir : - os que possuem grande domínio dos recursos de estruturação textual, com produção de textos que atendem às exigências de clareza e coerência, com progressão temática bem direcionada e utilização produtiva de operadores argumentativos ; - os que apresentam um domínio regular dos recursos de estruturação textual, com desvios na organização do texto e ausência de paragrafação, além de uso insuficiente de operadores argumentativos e processos de encaixamento sintático ; - os que apresentam baixo domínio dos recursos de estruturação textual, com muitos desvios que retomam hábitos da oralidade e que comprometem a coerência do texto : ocorrência de sequência justaposta de ideias sem encaixamentos sintáticos ; redução de estruturas subordinadas, ao lado do aumento na frequência de estruturas coordenadas e absolutas ; redução no uso de conectores para expressar relações lógicas essenciais à construção do texto ; emprego equivocado de operadores que não estabelecem relações lógicas coerentes entre ideias do texto ; fragmentação frasal que compromete a estrutura lógico-gramatical.
Uma estrutura que vem se ampliando na escrita dos participantes, refletindo hábitos da oralidade, é a duplicação de um dos elementos da frase, ou seja, a retomada anafórica de um elemento topicalizado, como por exemplo :"Elas precisam saber o verdadeiro papel DELAS na sociedade “, A desigualdade ELA vem de uma cultura machista... ” ; “ Homem e mulher, ELES são iguais ”.
Também em relação ao emprego dos sinais de pontuação, observaram -se distintos tipos de desempenho : participantes que dominam a funcionalidade da pontuação para a construção de textos fluentes e coesos, por um lado, e um grande número de participantes que apresentam precariedade em sua utilização, chegando ao ponto de elaborar parágrafos sem qualquer marca interna de pontuação para separar os períodos.
São os seguintes os tipos de desvios encontrados : - emprego de vírgula para separar o sujeito e o predicado : “ A mulher, sofre violência ” ( no lugar de “ A mulher sofre violência ” ) ;"O correto seria, ter um canal de denúncia ” ( no lugar de “ o correto seria ter um canal de denúncia ” ) ; - ocorrência de apenas uma das vírgulas para separar uma palavra, uma expressão ou uma oração encaixada : “ O homem, por sua vez não aceita que uma mulher seja sua chefa ” ( no lugar de “ O homem, por sua vez, não aceita que uma mulher seja sua chefa ” ) ; - emprego de ponto no lugar de vírgula para separar orações de um mesmo período : “ A todo custo as mulheres são atacadas na política.
Sexualmente também. ” ( no lugar de “ A todo custo as mulheres são atacadas na política, sexualmente também. ” ) ; - ausência de vírgula ou conjunção para separar elementos de uma enumeração :"... conquistas como voto equidade salarial direito de estudar e reconhecimento..."( no lugar de""... conquistas como voto, equidade salarial, direito de estudar e reconhecimento...") ; - ausência de vírgula para separar oração adjetiva explicativa ou utilização inadequada para separar oração adjetiva restritiva : “ As mulheres, que são iguais aos homens devem ser igualmente protegidas por lei."( no lugar de “ As mulheres, que são iguais aos homens, devem ser igualmente protegidas por lei.") ; - emprego indevido do sinal de dois pontos : “ Mulheres são julgadas : pelas simples roupas que usam. ” ( no lugar de “ Mulheres são julgadas pelas simples roupas que usam. ” ) ; - ausência de pronome no processo de referenciação : “ julga a mulher ou até agride ” ( no lugar de “ julga a mulher ou até a agride ” ” ) ; - emprego de vírgula em lugar do ponto-final : “ O patriarcado estrutural é formado por elementos históricos e religiosos, há pouquíssimo tempo, as mulheres passaram a votar. ” ( no lugar de “ O patriarcado estrutural é formado por elementos históricos e religiosos.
Há pouquíssimo tempo, as mulheres passaram a votar. ” ).
III - Aspectos Morfossintáticos e Vocabulares : o domínio das convenções morfossintáticas apresentou resultados que evidenciam, claramente, alguns aspectos problemáticos no desempenho dos participantes.
Assim como na análise dos desvios textuais, pode -se observar uma distinção entre três grupos de participantes : o grupo que revela domínio das convenções morfossintáticas ; o grupo que apresenta um domínio regular, com desvios pontuais ; e o grupo que revela domínio precário dessas convenções.
Pode -se destacar a regência, tanto verbal quanto nominal, como o aspecto com maior número de desvios : a ) Regência nominal e verbal Na maioria dos textos, destacam -se os desvios relacionados à noção de"regência", ou seja, à relação de subordinação que se estabelece entre as palavras por meio, ou não, de preposições, refletindo uma tendência na modalidade escrita da língua.
O desvio mais frequente foi a falta do sinal indicativo da crase, em contextos em que seria obrigatória, ou, ao contrário, um uso indevido em contextos em que o princípio não se aplica.
Os dois tipos de desvios revelam que o usuário não tem consciência de que, sob a forma do termo “ a ”, se pressupõe a presença de uma preposição “ a ”, exigida pela regência do termo anterior. tipos de desvios revelam que o usuário não tem consciência de que, sob a forma do termo “ a ”, se pressupõe a presença de uma preposição “ a ”, exigida pela regência do termo anterior.
Exemplos : “ diferentes tipos de agressão a mulher ” ( no lugar de “ diferentes tipos de agressão à mulher ” ) ; “ Desenvolverem práticas econômicas e sociais as novas gerações"( no lugar de “ Desenvolverem práticas econômicas e sociais às novas gerações") ; “ Além de proteger à mulher... ” ( no lugar de “ Além de proteger a mulher... ” ).
Outro desvio relacionado à regência verbal e à nominal, encontrado frequentemente nas questões, é a ausência de preposição antes de pronome relativo, refletindo um processo generalizado na modalidade oral da língua, em situações de registro informal.
Exemplo :"outra forma é oferecer apoio psicológico quem sofreu qualquer tipo de abuso..."( no lugar de :"outra forma é oferecer apoio psicológico a quem sofreu qualquer tipo de abuso...") b ) Concordância nominal e verbal Os desvios relativos à aplicação das regras de concordância verbal e nominal correspondem a fatos que têm sido amplamente estudados na modalidade oral da Língua Portuguesa.
Quanto à concordância de número, observou -se ausência de marca ( com sujeito anteposto ou posposto ), como por exemplo, a ausência de acento circunflexo na forma plural do presente do indicativo do verbo “ ter ”, considerada como um desvio na concordância verbal e não na acentuação gráfica.
Exemplo :"as mulheres ainda tem como obrigação"( no lugar de"as mulheres têm como obrigação") ; “ as leis em vigor tem que ser cumpridas ” ( no lugar de “ as leis em vigor têm que ser cumpridas ” ).
Outro desvio na aplicação da regra da concordância de número foi o aparecimento da marca de plural em verbos ou adjetivos relacionados a núcleos substantivos no singular, devido à presença de modificadores no plural, evidenciando um processo de hipercorreção.
Esse processo tem se ampliado nas situações de oralidade, caracterizando uma tendência de variação da Língua Portuguesa.
Exemplos :"Mudança nas leis que aumentem a pena do agressor"( no lugar de"Mudança nas leis que aumente a pena do agressor") ;"o aumento dos feminicídios devem ser combatidos ” ( no lugar de"o aumento dos feminicídios deve ser combatidos ” ).
Quanto à concordância de gênero, vários casos foram observados, normalmente no âmbito de sintagmas nominais longos, em que o adjetivo está afastado do substantivo, conforme apontam as pesquisas sobre a modalidade oral.
Exemplo : “ E as leis sobre as mulheres que sofrem violência ser mais rigorosos ” ( no lugar de “ E as leis sobre as mulheres que sofrem violência ser mais rigorosas ” ) ; Registro da forma verbal “ é ” sem acento, provocando ambiguidade na leitura pela confluência com a conjunção aditiva"e".
Exemplo : “ A desvantagem entre salários, poderes e entre a violência física e moral e crescente por demais... ” ( no lugar de “ A desvantagem entre salários, poderes e entre a violência física e moral é crescente por demais ”. c ) Emprego do pronome relativo “ onde ” como relativo universal, típica marca da oralidade, em que esse pronome não é mais utilizado apenas em sua função locativa"lugar onde se realiza a ação".
Exemplo :"Existe a desigualdade de gênero, onde homem ganha mais do que as mulheres"; “ há falta de informação e conhecimento onde ” ; “ O patriarcado está estruturado em nossa sociedade há séculos, onde o sexo masculino... ” d ) Falta de artigo definido antes de substantivo, comprometendo o paralelismo em enumerações Exemplo :"seja em oportunidade de emprego, salário"( no lugar de"seja em oportunidade de emprego, de salário"). e ) Desvios relacionados à flexão verbal Exemplos : emprego indevido do tempo ou do modo verbal : “ promovido ” por “ promovendo ” ; eliminação de marca de infinitivo : “ colocar ” por “ coloca ”, “ ta ” por “ estar ”, “ está ” por “ estar ”, “ esta ” por “ estar ” ; “ estás ” por “ estas ” ; utilização de construção de “ estar ” + gerúndio como marca de oralidade : “ a sociedade cada vez mais tá evoluindo ” ; segmentação indevida da desinência de número e pessoa de verbos : “ denunciar -mos ”.
Assim como nas demais competências avaliadas, há diferenças entre os participantes em relação à seleção e ao emprego do vocabulário.
Foram encontrados textos com alto padrão de domínio vocabular, chegando até ao emprego de vocabulário técnico das áreas do saber abordadas nas questões.
Por outro lado, na maioria dos textos, foram observados vários tipos de inadequação na seleção e na utilização do vocabulário, sendo que o principal deles foi a excessiva repetição de certas palavras, revelando limitação de repertório vocabular.
Além disso, observou -se a ocorrência de expressões da oralidade, revelando falta de domínio de vocabulário mais abstrato e de maior complexidade, essencial ao desenvolvimento do texto de base dissertativa.
Exemplos : seleção vocabular indevida : “ Em pleno século 21 é de emplena importância... ” ; emprego inadequado do verbo “ ter ” no lugar de “ haver ” : “ Durante a história, teve muito machismo ” ; imprecisão vocabular : “ Mas a cabeça, o corpo, há mesma igualdade para desenvolvermos igualmente as coisas ”, “ Em alguns anos atrás, as coisas eram mais complicadas nas questões de desigualdade no trabalho. ” A imprecisão vocabular, como o uso da palavra “ coisa ” em vários textos, demonstra o domínio precário da norma-padrão da Língua Portuguesa, com presença constante de traços de oralidade. 7.3.1.8 CONSIDERAÇÕES FINAIS SOBRE A CORREÇÃO DE LÍNGUA PORTUGUESA DAS RESPOSTAS ÀS QUESTÕES DISCURSIVAS DO COMPONENTE DE FORMAÇÃO GERAL O conjunto de textos analisados no Enade 2022 permite traçar um perfil bem detalhado do desempenho linguístico do estudante brasileiro de graduação. 7.3.2 COMPONENTE ESPECÍFICO A Tabela 7.17 apresenta dados sobre Questões Discursivas do Componente Específico.
A partir dos dados registrados na Tabela 7.17, observa -se que a média referente ao Componente de Conhecimento Específico foi mais alta do que para as Questões Discursivas do Componente de Formação Geral.
No Componente Específico, a média foi 41,2.
A maior média deste componente foi obtida pelos estudantes da região Sul ( 45,9 ), e a menor, pelos da região Norte ( 35,3 ).
As demais regiões obtiveram média igual a 41,5 ( Centro-Oeste ), 39,6 ( Nordeste ) e 42,5 ( Sudeste ).
Quanto à variabilidade das notas, o desvio-padrão de todo o Brasil foi 23,3.
O maior desvio-padrão foi encontrado na região Centro-Oeste ( 24,5 ), e o menor, na região Norte ( 22,7 ).
A maior nota máxima, 100,0, foi obtida por, pelo menos, um estudante na região Centro-Oeste.
Todas as regiões tiveram nota mínima zero.
A mediana do Brasil e da região Nordeste foi igual a 43,3 ; nas regiões Norte ( 36,7 ), Sul ( 50,0 ), Sudeste foi 46,7, e na região Centro-Oeste foi 40,0.
No Gráfico 7.10, representa -se a distribuição das notas nas Questões Discursivas no Componente Específico.
À luz dos dados do Gráfico 7.10, verifica -se que a moda da distribuição das notas nas Questões Discursivas no Componente Específico ocorre no intervalo ( 50 ; 60 ], com 17 \% do total de participantes.
A análise de cada uma destas questões será feita a seguir. 7.3.2.1 ANÁLISE DA QUESTÃO DISCURSIVA 3 DO COMPONENTE ESPECÍFICO Os resultados da Questão Discursiva 3 do Componente Específico são descritos na Tabela 7.18 e comentados na sequência.
Como mostrado na Tabela 7.18, na Questão 3, a média dos estudantes de todo o Brasil foi 23,7.
A menor média na questão, em evidência, foi obtida pelos estudantes da região Norte ( 19,1 ), enquanto a maior média foi obtida na região Sul ( 27,2 ).
Quanto à variabilidade das notas, o desvio-padrão de todo o Brasil foi 25,1.
O maior desvio-padrão foi obtido na região Centro-Oeste ( 28,4 ), enquanto o menor foi obtido na região Norte ( 21,4 ).
A nota máxima, 100,0 pontos, foi alcançada por, pelo menos, um estudante em cada região.
A mediana do Brasil foi 20,0, para maioria das regiões ; exceto nas regiões Nordeste e Norte que foram, respectivamente 10,0.
A nota mínima zero foi a mesma em todas as regiões do Brasil.
Destaca -se, com base nos dados do Gráfico 7.11, que 46,0 \% dos estudantes deixaram a resposta à questão em branco.
A moda dessa distribuição ocorre no intervalo ( 0 ; 10 ], seguido do intervalo ( 10 ; 20 ], com 46 \% e 18 \% do total de participantes, respectivamente. 7.3.2.2 COMENTÁRIOS SOBRE AS RESPOSTAS À QUESTÃO DISCURSIVA 3 A Questão 3 abordou um tema de grande relevância para a formação de nível superior em Gastronomia, trazida à luz, sobretudo, pela pandemia da Covid-19, que é o papel do manipulador de alimentos para a garantia da qualidade higiênico-sanitária dos alimentos e para a manutenção da saúde dos comensais.
O tema pertence às Diretrizes Curriculares Nacionais do Curso, sendo desenvolvido ao longo de toda a formação dos egressos, tanto nos componentes teóricos quanto práticos de laboratório.
Com isso, esperava -se do estudante o domínio sobre os conceitos de limpeza e sanitização definidos pela Agência Nacional de Vigilância Sanitária ( Anvisa ).
Observou -se, no entanto, dificuldades de os estudantes definirem com clareza as diferenças existentes entre um conceito e outro, o que gerou impacto sobre a resposta do melhor momento em realizar a sanitização, conforme solicitado no item “ b ” da Questão.
Ainda em relação a esse item, era esperada mais facilidade em indicar 2 ( dois ) exemplos de sanitizantes uma vez que fazem uso de álcool 70 \% e solução de cloro com regularidade nos laboratórios de gastronomia.
Boa parte dos estudantes respondeu de forma diferente ao que constava no padrão de resposta.
No item “ a ”, um comportamento muito recorrente foi que os estudantes usavam os temos “ limpeza ” e “ sanitização ” de modo equivalente, demonstrando desconhecimento sobre a diferença existente entre os dois conceitos.
Muitas respostas não abordavam nenhum conhecimento técnico sobre as normativas da ANVISA para a questão.
Outros argumentavam que, ora a limpeza, ora a sanitização, deviam ser observadas para evitar a contaminação cruzada dos alimentos ou causar prejuízos à saúde dos comensais.
Poucos estudantes demonstraram clareza sobre os conceitos e responderam os itens “ a ” e “ b ” conforme o padrão de resposta.
Sobre os exemplos de sanitizantes, as citações mais usuais foram “ água sanitária ( e suas variações ) ” e “ álcool 70 \% ”, embora tenham sido observados muitos esquecimentos em relação ao percentual mínimo do álcool e exemplos de marcas comerciais de água sanitária ( como “ Qboa ” ).
Com relação à linguagem utilizada pelos estudantes, as respostas demonstram dificuldades de responder as questões de forma objetiva, separadas em tópicos, conforme indicado no enunciado.
Alguns estudantes definiram os conceitos com base em termos técnicos, enquanto outros usaram expressões que demonstram noções das boas práticas, porém nem sempre estabelecidas no padrão de resposta, como “ limpar a bancada ” ou “ trocar as tábuas ”, entre outros.
Com relação ao item “ a ” da Questão 3, alguns estudantes responderam que “ a limpeza e a sanitização devem ser feitas a fim de reduzir a contaminação do ambiente, garantindo a saúde do consumidor ”, sem distinguir entre um conceito e outro.
Alguns estudantes descreveram que o momento ideal para realizar a sanitização é quando os ingredientes são recebidos na cozinha ou antes de dar início ao preparo.
Uma resposta bastante usual foi relacionada aos cuidados do manipulador com a lavagem das mãos, uso de toucas e uniforme.
Também foi comum encontrar respostas relacionadas à dedetização do ambiente e controle de vetores e pragas, informações que não foram demandadas na Questão 3.
Alguns estudantes demonstraram um bom conhecimento sobre os conceitos, indicando, inclusive, a resolução da Anvisa que trata do tema como exemplo.
Estes, não apenas indicaram os conceitos conforme o padrão de resposta, mas souberam indicar o melhor momento para realizar a sanitização, como demandado no item “ b ” e indicaram dois exemplos distintos de sanitizantes.
Ainda com relação ao item “ b ”, alguns citaram corretamente que o melhor momento para realizar a sanitização é após a limpeza, demonstrando falta de clareza na Questão.
Outra resposta encontrada com frequência foi a indicação de dois exemplos iguais de sanitizantes, “ água sanitária ou cloro diluído em água ”.
Sobre os exemplos, foi comum encontrar apenas a citação de álcool, sem indicação do percentual mínimo de 70 \%.
Em relação ao padrão de resposta do item “ a ”, a distribuição de pontos segundo a régua de graduação foi essencial para que se pudesse aferir alguma pontuação aos estudantes.
Com relação ao item “ b ”, muitos não conseguiram responder adequadamente a primeira parte, mas outros citaram pelo menos um exemplo de sanitizantes.
A maioria dos estudantes respondeu à Questão 3, no entanto, o desempenho foi abaixo do esperado, pois o enunciado pedia a definição de conceitos segundo a Anvisa.
As respostas demonstravam que os estudantes tinham a noção sobre a importância das Boas Práticas na Manipulação de Alimentos, porém sem aprofundamento em relação aos conceitos de limpeza e sanitização.
Além disso, nota -se que o maior enfoque se deu sobre a higienização dos alimentos ou controle do manipulador, como lavagem das mãos ou uso de uniforme, demonstrando menor conhecimento em relação à limpeza e sanitização de equipamentos, utensílios, ou mesmo, em relação ao espaço físico de trabalho.
Muitos estudantes citaram pelo menos um exemplo de sanitizante corretamente.
A maior dificuldade foi encontrada no item “ a ”, que exigia conhecimento das normas da Anvisa e, além disso, a organização do texto para responder separadamente o conceito de “ limpeza ” e “ sanitização ”.
Em razão do desconhecimento, muitos responderam como sendo um conceito único.
Por outro lado, os estudantes que responderam adequadamente sabiam citar, inclusive, a normativa da Anvisa que trata do tema.
As respostas medianas abordavam ao menos um conceito adequadamente e citavam um ou dois exemplos corretos de sanitizantes.
Já as respostas consideradas boas foram aquelas em que os estudantes respondiam sobre os dois conceitos exigidos de maneira mais completa, indicaram o melhor momento para a sanitização, porém esqueceram algum exemplo de sanitizante ou não responderam integralmente como o padrão de resposta exigia.
As respostas indicam que a maioria dos estudantes tem noção da importância das boas práticas e seu papel como manipuladores de alimentos, porém demostram pouco domínio sobre os conceitos técnicos estabelecidos pela Anvisa.
Isso pode indicar um distanciamento entre a teoria e a prática de laboratório ou mesmo profissional.
Pela estrutura das respostas, os estudantes demonstram dificuldade em responder os itens exigidos no enunciado com clareza e organização. 7.3.2.3 ANÁLISE DA QUESTÃO DISCURSIVA 4 DO COMPONENTE ESPECÍFICO Na Tabela 7.19, constam as informações relativas à Questão 4 do conjunto de Questões do Componente Específico.
Como se observa, a partir dos dados da Tabela 7.19, o desempenho dos estudantes nessa questão foi superior ao da Questão 3 e superior ao da Questão 5.
A média geral do Brasil foi 58,7, sendo a maior média registrada na região Sul ( 63,0 ), e a menor, na região Norte ( 51,9 ).
Quanto à variabilidade das notas, o desvio-padrão geral do Brasil foi 30,8, sendo o maior desvio-padrão registrado na região Norte ( 32,0 ), e o menor, na região Nordeste ( 29,7 ).
A nota máxima ( 100,0 ) foi atingida por, pelo menos, um concluinte em todas as regiões.
A mediana em todo o Brasil foi 70,0, o mesmo valor foi obtido em todas as regiões, exceto nas regiões Norte e Nordeste ( 60,0 ).
Foi zero a nota mínima de todas as regiões.
No Gráfico 7.12, está representada a distribuição de notas da Questão Discursiva 4 do Componente Específico.
A distribuição mostrada no Gráfico 7.12 tem moda nos intervalos ( 60 ; 70 ] e [ 80 ; 90 ] com pouco mais de 15 \% do total.
Chama a atenção que, nessa questão, mais de 15 \% obtiveram nota zero. 7.3.2.4 COMENTÁRIOS SOBRE AS RESPOSTAS À QUESTÃO DISCURSIVA 4 A Questão 4 da prova específica da área de Gastronomia diz respeito aos cortes clássicos e métodos e técnicas de cocção, conhecimentos, competências e habilidades básicas para a prática da gastronomia, especialmente para aqueles que irão trabalhar em uma cozinha profissional.
O comando da questão solicitava a identificação e explicação de três cortes e duas técnicas/métodos.
O padrão de respostas aceita como correto seis cortes e cinco técnicas/métodos, sendo um ponto para cada corte identificado ( máximo 3 pontos ), 1 ponto para cada descrição correspondente ao corte identificado ( máximo 3 pontos ), 1 ponto para cada técnica/método identificado ( máximo 2 pontos ) e 1 ponto para cada descrição a técnicas/métodos identificados ( máximo 2 pontos ).
Ao que se refere aos cortes, item “ a ” da questão, o corte concassé se destacou por ser identificado pelos estudantes recorrentemente como técnica/método.
Referente ao corte tourné ( ou torneado ) as respostas foram em grande maioria diferentes da resposta padrão.
Em relação às técnicas/métodos, as palavras grelhar e saltear por vezes foram confundidas uma com a outra e também tiveram as respostas muitas vezes diferentes, incompletas se comparadas à resposta padrão.
De modo geral, as respostas estavam contidas no espaço destinado a elas, seja em forma de texto corrido ou separadas nos itens “ a ” e “ b ”.
Observou -se considerável quantidade de erros de grafia nas palavras dos cortes e técnicas/métodos mesmo elas estando contidas no enunciado da questão.
Algumas respostas explicaram os cortes e métodos/técnicas minuciosamente extrapolando a resposta padrão.
Outras, no entanto foram demasiadamente econômicas sendo desconsideradas como correta na avaliação, como por exemplo, “ grelhar : fazer uso da grelha ” ou “ saltera : saltear na saltese ”.
Boa parte dos estudantes respondeu à questão 4.
As respostas demonstraram que os estudantes identificam os cortes e técnicas/métodos, mas não o explicam completamente.
Percebe -se que a maioria tem uma noção das características dos cortes e técnica/método.
Não houve desrespeito aos direitos humanos.
As respostas consideradas fracas foram aquelas nas quais todos os itens foram respondidos erroneamente ou parcialmente.
As respostas medianas indicavam os cortes e técnicas/métodos sem explicá -las.
As respostas consideradas boas foram aquelas em que os estudantes respondiam sobre o indicaram os cortes e técnicas/métodos e apresentavam no mínimo duas das características apresentadas na resposta padrão.
A indicação dos cortes e métodos/técnicas foram apresentados em grande parte das respostas.
A explicação por vezes não foi apresentada.
Uma característica é que ou o estudante sabia a explicação dos cortes ou dos métodos/técnicas.
Ficou evidente que os estudantes têm conhecimento raso sobre o tema, precisando de um aprofundamento.
Sendo o tema relativo à área prática, fica clara a necessidade de mais horas de práticas de laboratório para fixação do conteúdo. 7.3.2.5 ANÁLISE DA QUESTÃO DISCURSIVA 5 DO COMPONENTE ESPECÍFICO Na Tabela 7.20, constam as informações relativas à Questão 5 do conjunto do Componente Específico.
Como revelam os dados da Tabela 7.20, o desempenho dos estudantes de todo o Brasil na Questão Discursiva 5 do Componente Específico foi inferior ao das Questões Discursivas 4, desse componente.
A nota média dos estudantes de todo o Brasil foi 41,1,0.
A maior média foi registrada na região Sul ( 47,6 ), enquanto a menor média foi registrada na região Centro-Oeste ( 38,9 ).
Quanto à variabilidade das notas, o desvio-padrão dos estudantes do Brasil foi 31,1.
Enquanto o maior desvio foi encontrado na região Nordeste ( 31,6 ), o menor foi encontrado na região Sul ( 30,2 ).
A mediana para o Brasil foi 50,0.
Para o conjunto de estudantes de Tecnologia em Gastronomia do Brasil, a nota máxima foi 100,0, obtida por, pelo menos, um estudante em cada região.
A nota mínima do Brasil foi zero, o mesmo valor encontrado em todas as regiões.
No Gráfico 7.13, é representada a distribuição das notas da Questão Discursiva 5 do Componente Específico.
A distribuição mostrada no Gráfico 7.13 tem moda no intervalo ( 60 ; 70 ].
Nessa questão, 6,9 \% obtiveram nota zero e 14,8 \% deixaram a questão em branco.
A soma das frequências destes dois intervalos \textbf{ultrapassa} 1/5 dos respondentes ( 21,7 \% ). 7.3.2.6 COMENTÁRIOS SOBRE AS RESPOSTAS À QUESTÃO DISCURSIVA 5 As respostas, de um modo geral, contemplam o que o item pede no enunciado.
No item “ a ”, o aluno necessitou possuir o conhecimento do conceito de custo variável, que ele poderia aplicar dentro de uma empresa, ao exemplificar, pelo menos, um custo.
No item “ b ”, o aluno necessitou possuir o conhecimento do conceito de custo fixo, que ele poderia aplicar dentro de uma empresa, exemplificando, pelo menos, um custo.
Vale ressaltar que o aluno deveria falar a respeito da diferença, já que os conceitos são muito próximos e, às vezes, podem ser confundidos.
Já na questão “ c ”, o aluno deveria ter o conhecimento de como é produzido e precificado um alimento para a venda.
Essa questão envolve uma multidisciplinaridade, na qual é necessário que o aluno conheça o fator de correção na produção, os custos de produção, a relação entre fator de correção e os custos de produção e indique a existência de perdas no processo e a utilização do fator de correção para o cálculo correto dos custos de produção dos alimentos.
Ela exige que o aluno tenha o entendimento de vários fatores que incidem desde a compra, produção e venda de um alimento, fechando, assim, a cadeia de produção.
A Questão 5 exige que o aluno tenha conhecimento prático e teórico, já que os itens pedidos no enunciado têm grande grau de confusão, já que os conceitos se assemelham bastante.
Com relação ao item “ a ”, alguns estudantes responderam que “ o que era custo variável ”, sem distinguir o conceito de “ custo fixo ”, já que era o conhecimento solicitado no item “ b ”.
Vários estudantes faziam grande confusão entre os dois termos, ora podiam acertar a descrição e errar o exemplo, ora podiam trocar a descrição dos dois, colocando a explicação de custo fixo no lugar de custo variável.
Uma porcentagem dos alunos também errou a exemplificação, podendo confundir a aplicação dos dois custos, sendo assim penalizado de um a dois pontos.
Já no terceiro quesito da questão, grande parte dos alunos entendiam que era muito importante o Fator de Correção, porém uma grande parte não sabia o que era e por que era utilizada essa ferramenta dentro da gastronomia.
Em várias respostas, foi apenas afirmado que ela era importante, sem a justificativa, atribuindo assim apenas um ponto para essa afirmativa.
Mas nesse ponto, entende -se que o aluno necessitava de um completo entendimento do funcionamento de um estabelecimento comercial, já que ele deveria entender as variações diárias, a necessidade de uma ficha técnica bem elaborada para assim fazer compras assertivas e um bom cálculo para vender seus produtos nesses estabelecimentos gastronômicos.
A Questão 5 foi a que teve maior abstenção, as colegas que fizeram as correções das Questões 3 e 4, relataram -me que houve uma maior assiduidade nas respostas, creio que essa abstenção, foi relacionada ao cansaço, já que era a última questão subjetiva.
Além disso, foram corrigidas algumas provas com : “ estou muito cansado ”, “ a prova foi muito longa ”, entre outras afirmações.
No entanto o aluno, para responder esses itens solicitados na Questão 5, precisava possuir o conhecimento teórico e a aplicação desse conteúdo na prática, logo em inúmeras respostas, o aluno apenas colocava os exemplos e não conseguia fazer a definição do item.
Em outras respostas, os alunos conseguiam fazer a definição, porém se confundiam com o exemplo.
Sendo um exemplo extremamente utilizado para definir custo variável ( “ energia elétrica ” e “ conta de Água ” ), faz -se necessário a \textbf{observação} de um entendimento lógico, já que é um item na planilha de custo que pode haver variação dependendo a produção e a utilização, mas é um item que deverá ser pago todo mês, e sem esses dois itens não seria possível o funcionamento do estabelecimento gastronômico, tornando -se um custo fixo ( ARAUJO, Adriana Maria Procópio de, ASSAF, Alexandre.
Introdução à Contabilidade.
São Paulo : Atlas, 2004 ).
No último item, os alunos sabiam da importância do fator de correção, porém, não sabiam a aplicabilidade da ferramenta, fazendo com que muitos não fechassem o componente, mas praticamente todos os alunos, responderam positivamente sobre a importância da ferramenta, sendo atribuído assim um ponto.
As respostas consideradas fracas foram aquelas nas quais todos os itens foram respondidos erroneamente ou apenas parcialmente.
As respostas medianas abordavam ao menos um conceito adequadamente e citavam um ou dois exemplos corretos em custo fixo ou custo variável ou a importância do fator de correção.
Já as respostas consideradas boas foram aquelas em que os estudantes respondiam sobre os dois conceitos de custo fixo e custo variável de maneira mais completa, indicando os exemplos corretos, falando da importância do fator de correção e sabendo a aplicabilidade do mesmo.
Era possível ver que os alunos conseguiram ver esses itens cobrados na questão, porém era claro que grande parte dos alunos não absorveu esse conteúdo, já que os alunos ora sabiam a definição e não sabiam o exemplo, ora sabiam o exemplo e não sabiam a definição, ora trocavam os exemplos.
Muitos deles sabiam pelo menos um item.
Essa questão faz -se necessária para a formação do Gastrônomo, já que é de suma importância a gestão de custos, gastos, vendas e lucro de um estabelecimento.
No custo fixo e variável, o profissional que for empreender precisa de saber bem definido quais são, já que, para uma boa saúde financeira, ele necessita entrar em ponto de equilíbrio ( conseguindo equilibrar venda e custos ), e nesse ponto o profissional pode administrar esses dois itens a fim de economizar e achar modos para que esses valores não sufoque a empresa.
Já no item “ c ”, o gastrônomo que irá trabalhar, sendo ele empresário ou funcionário, necessita desse conhecimento para uma boa administração da produção, do estoque e das compras de insumos.
Tanto que esse item traz a importância da ficha técnica, e nela necessita de informações técnicas para : produção, compra e cálculo de venda do item.
A Questão 5 tratava de um assunto relevante na formação do Gastrônomo, especialmente aqueles que irão trabalhar com empreendedores, na produção ou mesmo na gestão de empreendimentos gastronômicos. 7.3.2.7 CONSIDERAÇÕES FINAIS SOBRE AS QUESTÕES DISCURSIVAS DO COMPONENTE ESPECÍFICO Poucos estudantes demonstraram clareza sobre os conceitos, e foram observados muitos esquecimentos em relação a definições e/ou exemplos, ou ainda uma confusão entre dois termos.
Percebeu -se também que as ferramentas são bem exploradas dentro do campo teórico, já que grande maioria respondeu algo, porém, somente com a prática dentro de estabelecimentos funcionais, os alunos irão ter um completo conhecimento da teoria e da aplicabilidade.
Logo atribuo a esse fator o motivo de boa parte dos alunos não conseguirem desenvolver muito os conceitos.
Após essa análise, fica clara a necessidade de mais horas de práticas de laboratório para fixação do conteúdo.

% matched lemmas: brasileira, observação, ultrapassar
\end{textsample}
